\frfilename{mt283.tex}
\versiondate{31.3.13}

\def\chaptername{Analisis Fourier}
\def\sectionname{Transformasi Fourier I}
\copyrightdate{1994}

\newsection{283}

Sekarang saya beralih ke teori transformasi Fourier pada $\Bbb R$. Dalam
bagian pertama dari dua bagian mengenai topik ini, saya menyajikan bagian teori
dasar yang dapat ditangani dengan metode dari bagian sebelumnya mengenai deret
Fourier. Namun, saya tidak menemukan cara untuk memahami teori ini tanpa
memperkenalkan sebagian teori distribusi L.Schwartz, yang saya sajikan dalam
\S284. Seperti dalam \S282, pembahasan ini tentu saja baru merupakan suatu
permulaan.

Seluruh teori ini juga dapat dikembangkan pada $\BbbR^r$. Namun saya
menyerahkan perluasan tersebut kepada latihan, karena hanya ada sedikit
gagasan baru, rumus-rumusnya jauh lebih rumit, dan setidaknya dalam jilid ini
saya tidak memerlukan versi multidimensi dari teorema-teorema tersebut,
meskipun sebagian gagasan yang sama akan muncul dalam bentuk multidimensi
dalam \S285.

\leader{283A}{Definisi} Misalkan $f$ adalah fungsi bernilai kompleks yang
dapat diintegralkan pada $\Bbb R$.


\spheader 283Aa {\bf Transformasi Fourier} dari $f$ adalah fungsi
  $\varhatf:\Bbb R\to\Bbb C$ yang didefinisikan oleh

\Centerline{$\varhatf(y)=\Bover1{\sqrt{2\pi}}\int_{-\infty}^{\infty}
e^{-iyx}f(x)dx$}

\noindent untuk setiap $y\in\Bbb R$.   \cmmnt{(Tentu saja integral
selalu ditentukan karena $x\mapsto e^{-iyx}$ terbatas dan kontinu,
oleh karena itu terukur.)}

\spheader 283Ab  {\bf transformasi Fourier invers} dari $f$ adalah
fungsi $\varcheckf:\Bbb R\to\Bbb C$ yang ditentukan oleh

\Centerline{$\varcheckf(y)=\Bover1{\sqrt{2\pi}}\int_{-\infty}^{\infty}
e^{iyx}f(x)dx$}

\noindent untuk setiap $y\in\Bbb R$.

\cmmnt{
\leader{283B}{Keterangan (a)} Ini adalah fitur yang agak menjengkelkan dari teori
dari transformasi Fourier -- menjengkelkan, yaitu,
untuk orang luar seperti saya -- bahwa sebenarnya
tidak ada definisi standar dari `transformasi Fourier'.
Definisi yang paling umum, menurut saya, adalah,

\Centerline{$\varhatf(y)=\Bover1{\sqrt{2\pi}}\int_{-\infty}^{\infty}
e^{\mp iyx}f(x)dx$,}

\Centerline{$\varhatf(y)=\int_{-\infty}^{\infty}
e^{\mp iyx}f(x)dx$,}

\Centerline{$\varhatf(y)=\int_{-\infty}^{\infty}
e^{\mp 2\pi iyx}f(x)dx$,}

\noindent sesuai dengan transformasi Fourier invers

\Centerline{$\varcheckf(y)=\Bover1{\sqrt{2\pi}}\int_{-\infty}^{\infty}
e^{\pm iyx}f(x)dx$,}

\Centerline{$\varcheckf(y)=\Bover1{{2\pi}}\int_{-\infty}^{\infty}
e^{\pm iyx}f(x)dx$,}

\Centerline{$\varcheckf(y)=\int_{-\infty}^{\infty}
e^{\pm 2\pi iyx}f(x)dx$.}

\noindent Saya serahkan kepada pembaca untuk memeriksa bahwa seluruh teori dapat
dijalankan dengan sembarang di antara keenam pasangan ini, serta menyelidiki
kemungkinan-kemungkinan lain (lihat 283Xa-283Xb di bawah).

\spheader 283Bb Ungkapan `transformasi Fourier' dan `transformasi Fourier
invers' menunjukkan bahwa $(\varhatf)\varspcheck$ semestinya sama dengan
$f$, setidaknya dalam beberapa keadaan. Hal ini memang benar, tetapi kelas
fungsi $f$ yang memenuhinya secara harfiah agak terbatas; kita masih harus
menunggu sebelum menyelidikinya.

\spheader 283Bc Betapa pun konstanta-konstanta diatur seperti pada (a),
kita tidak dapat membuat $\varhatf$ dan $\varcheckf$ benar-benar sama.
Namun, dengan definisi yang saya pilih, kita mempunyai
$\varcheckf(y)=\varhatf(-y)$ untuk setiap $y$, sehingga $\varcheckf$ dan
$\varhatf$ pada dasarnya memiliki semua sifat yang sama yang kita perlukan
di sini; khususnya, segala sesuatu dalam proposisi berikutnya akan valid
dengan $\varspcheck$ di tempat $\varsphat$, jika Anda mengubah tanda di
    titik-titik yang sesuai pada bagian (c), (h), dan (i).
}%end of comment

\leader{283C}{Proposisi} Misalkan $f$ dan $g$ adalah fungsi bernilai kompleks
yang dapat diintegralkan pada $\Bbb R$.

(a) $(f+g)\varsphat=\varhatf+\varhat{g}$.

(b) $(cf)\varsphat=c\varhatf$ untuk setiap $c\in\Bbb C$.

(c) Jika $c\in\Bbb R$ dan $h(x)=f(x+c)$ kapan pun ruas ini
terdefinisi, maka $\varhat{h}(y)=e^{icy}\varhatf(y)$ untuk setiap $y\in\Bbb R$.

(d) Jika $c\in\Bbb R$ dan $h(x)=e^{icx}f(x)$ untuk setiap
$x\in\dom f$, lalu $\varhat{h}(y)=\varhatf(y-c)$ untuk setiap
$y\in\Bbb R$.

(e) Jika $c>0$ dan $h(x)=f(cx)$ kapan pun ruas ini terdefinisi,
lalu $\varhat{h}(y)=\Bover1c\varhatf(\Bover{y}{c})$
untuk setiap $y\in\Bbb R$.

(f) $\varhatf:\Bbb R\to\Bbb C$ kontinu.

(g) $\lim_{y\to\infty}\varhatf(y)=\lim_{y\to-\infty}\varhatf(y)=0$.

(h) Jika $\int_{-\infty}^{\infty}|xf(x)|dx<\infty$, maka $\varhatf$ adalah
terdiferensiasi, dan turunannya adalah

\Centerline{$\varhatf\vthsp'(y)=-\Bover{i}{\sqrt{2\pi}}
\int_{-\infty}^{\infty}e^{-iyx}xf(x)dx$}

\noindent untuk setiap $y\in\Bbb R$.

(i) Jika $f$ kontinu mutlak pada setiap interval terbatas dan $f'$
dapat diintegralkan, maka $(f')\varsphat(y)=iy\varhatf(y)$ untuk setiap $y\in\Bbb
R$.

\proof{{\bf (a)} dan {\bf (b)} adalah sepele, dan {\bf (c)}, {\bf (d)} dan
{\bf (e)} adalah substitusi dasar.

\medskip

{\bf (f)} Jika $\sequencen{y_n}$ adalah sembarang barisan konvergen di $\Bbb R$
dengan limit $y$, maka

$$\eqalign{\varhatf(y)
&=\Bover1{\sqrt{2\pi}}
\int_{-\infty}^{\infty}\lim_{n\to\infty}e^{-iy_nx}f(x)dx\cr
&=\lim_{n\to\infty}\Bover1{\sqrt{2\pi}}
\int_{-\infty}^{\infty}e^{-iy_nx}f(x)dx
=\lim_{n\to\infty}\varhatf(y_n)\cr}$$

\noindent oleh Teorema Konvergensi Dominasi Lebesgue, karena
$|e^{-iy_nx}f(x)|\le|f(x)|$ untuk setiap $n\in\Bbb N$ dan $x\in\dom f$.   Karena
$\sequencen{y_n}$ bersifat arbitrer, $\varhatf$ bersifat kontinu.

\medskip

{\bf (g)} Ini hanya lema Riemann--Lebesgue (282E).

\medskip

{\bf (h)} Intinya $|\pd{}{y}e^{-iyx}f(x)|=|xf(x)|$ kapan saja
$x\in\dom f$ dan $y\in\Bbb R$. Jadi oleh 123D

$$\eqalign{\varhatf\vthsp'(y)
&=\Bover1{\sqrt{2\pi}}\Bover{d}{dy}
     \int_{-\infty}^{\infty}e^{-iyx}f(x)dx
=\Bover1{\sqrt{2\pi}}\Bover{d}{dy}\int_{\dom f}
     e^{-iyx}f(x)dx\cr
&=\Bover1{\sqrt{2\pi}}\int_{\dom f}
     \Bover{\partial}{\partial y}e^{-iyx}f(x)dx
=\Bover1{\sqrt{2\pi}}\int_{-\infty}^{\infty}
     -ixe^{-iyx}f(x)dx\cr
&=-\Bover{i}{\sqrt{2\pi}}\int_{-\infty}^{\infty}
     xe^{-iyx}f(x)dx.\cr}$$

\medskip

{\bf (i)} Karena $f$ kontinu mutlak di setiap interval terbatas,

\Centerline{$f(x)=f(0)+\int_0^xf'$ untuk $x\ge 0$,\quad
$f(x)=f(0)-\int_x^0f'$ untuk $x\le 0$.}

\noindent Karena $f'$ dapat diintegralkan,

\Centerline{$\lim_{x\to\infty}f(x)=f(0)+\int_0^{\infty}f'$,\quad
$\lim_{x\to-\infty}f(x)=f(0)-\int_{-\infty}^0f'$}

\noindent keduanya ada.   Karena $f$ juga dapat diintegralkan, kedua batas
harus nol.   Sekarang kita dapat melakukan integrasi parsial (225F) untuk melihat bahwa

$$\eqalign{(f')\varsphat(y)
&=\Bover1{\sqrt{2\pi}}\int_{-\infty}^{\infty}e^{-iyx}f'(x)dx
=\Bover1{\sqrt{2\pi}}\lim_{a\to\infty}\int_{-a}^{a}e^{-iyx}f'(x)dx\cr
&=\Bover1{\sqrt{2\pi}}\bigl(\lim_{a\to\infty}e^{-iya}f(a)
-\lim_{a\to-\infty}e^{-iya}f(a)\bigr)
+\Bover{iy}{\sqrt{2\pi}}\int_{-\infty}^{\infty}e^{-iyx}f(x)dx\cr
&=iy\varhatf(y).\cr}$$
}%end of proof of 283C

\leader{283D}{Lema} (a) $\lim_{a\to\infty}\int_{0}^a\bover{\sin
x}{x}dx=\bover{\pi}2$, $\lim_{a\to\infty}\int_{-a}^a\bover{\sin
x}{x}dx=\pi$.

(b) Terdapat $K<\infty$ sehingga $|\int_a^b\bover{\sin cx}{x}dx|\le
K$ kapan pun $a\le b$ dan $c\in\Bbb R$.

\proof{{\bf (a)(i)} Atur

\Centerline{$F(a)=\int_0^a\Bover{\sin x}{x}dx$ jika $a\ge 0$,\quad
$F(a)=-\int_{-a}^0\Bover{\sin x}{x}dx$ jika $a\le 0$,}

\noindent sehingga $F(a)=-F(-a)$ dan
$\int_a^b\bover{\sin x}{x}dx=F(b)-F(a)$ untuk semua $a\le b$.

Jika $0<a\le b$, maka menurut 224J

\Centerline{$|\int_a^b\Bover{\sin x}{x}dx|
\le(\Bover1b+\Bover1a-\Bover1b)\sup_{c\in[a,b]}
  |\int_a^c\sin x\,dx|
\le\Bover1a\sup_{c\in[a,b]}|\cos a-\cos c|
\le\Bover2a$.}

\noindent Khususnya, $|F(n)-F(m)|\le\bover2{m}$ jika $0<m\le n$ dalam
$\Bbb N$, dan $\sequencen{F(n)}$ adalah urutan Cauchy dengan batas
$\gamma$ katakan; sekarang

\Centerline{$|\gamma - F(a)|
=\lim_{n\to\infty}|F(n)-F(a)|\le\Bover2a$}

\noindent untuk setiap $a>0$, Jadi $\lim_{a\to\infty}F(a)=\gamma$.
Tentu saja kita juga memiliki

\Centerline{$\lim_{a\to\infty}\int_{-a}^a\Bover{\sin x}{x}dx
=\lim_{a\to\infty}(F(a)-F(-a))
=\lim_{a\to\infty}2F(a)=2\gamma$.}

\medskip

\quad{\bf (ii)} Jadi sekarang saya harus menghitung $\gamma$.    Untuk ini,
amati dulu

\Centerline{$2\gamma
=\lim_{a\to\infty}\int_{-\pi a}^{\pi a}\Bover{\sin x}{x}dx
=\lim_{a\to\infty}\int_{-\pi}^{\pi}\Bover{\sin at}{t}dt$}

\noindent (dengan mengganti $x=at$).   Berikutnya,

\Centerline{$\lim_{t\to 0}\Bover1t-\Bover1{2\sin{1\over 2}t}=\lim_{u\to
0}\Bover{\sin u-u}{2u\sin u}=0$,}

\noindent jadi

\Centerline{$\int_{-\pi}^{\pi}
  \bigl|\Bover1t-\Bover1{2\sin{1\over 2}t}\bigr|dt<\infty$,}

\noindent dan oleh lema Riemann-Lebesgue (282Fb)

\Centerline{$\lim_{a\to\infty}\int_{-\pi}^{\pi}
   \bigl(\Bover1t-\Bover1{2\sin{1\over 2}t}\bigr)\sin at\,dt=0$.}

\noindent Tapi kami tahu itu

\Centerline{$\int_{-\pi}^{\pi}
\Bover{\sin(n+{1\over 2})t\pushbottom{3.5pt}}{2\sin{1\over 2}t}dt
=\pi$}

\noindent untuk setiap $n$ (menggunakan 282Dc), jadi kita harus punya

$$\eqalign{\lim_{a\to\infty}\int_{-a}^a\Bover{\sin t}{t}dt
&=\lim_{a\to\infty}\int_{-\pi}^{\pi}\Bover{\sin at}{t}dt
=\lim_{a\to\infty}\int_{-\pi}^{\pi}
      \Bover{\sin at}{2\sin{1\over 2}t}dt\cr
&=\lim_{n\to\infty}\int_{-\pi}^{\pi}
    \Bover{\sin(n+{1\over 2})t\pushbottom{3.5pt}}{2\sin{1\over 2}t}dt
=\pi,\cr}$$

\noindent dan $\gamma=\pi/2$, seperti yang diklaim.

\medskip

{\bf (b)} Karena $F$ kontinu dan

\Centerline{$\lim_{a\to\infty}F(a)=\gamma=\Bover{\pi}2$, \quad
$\lim_{a\to-\infty}F(a)=-\gamma=-\Bover{\pi}2$,}

\noindent $F$ terbatas; katakanlah $|F(a)|\le K_1$ untuk semua $a\in\Bbb R$.
Ambil $K=2K_1$. Sekarang misalkan $a<b$ dan $c\in\Bbb R$. Jika $c>0$,
maka

\Centerline{$|\int_a^b\Bover{\sin cx}{x}dx|
=|\int_{ac}^{bc}\Bover{\sin t}{t}dt|
=|F(bc)-F(ac)|\le 2K_1=K$,}

\noindent menggantikan $x=t/c$.   Jika $c<0$, maka

\Centerline{$|\int_a^b\Bover{\sin cx}{x}dx|
=|-\int_a^b\Bover{\sin(-c)x}{x}dx|
\le K$;}

\noindent sementara jika $c=0$ maka

\Centerline{$|\int_a^b\Bover{\sin cx}{x}dx|
=0\le K$.}
}%end of proof of 283D

\leader{283E}{}\cmmnt{ Pekerjaan tersulit dari bagian ini terletak pada
`teorema inversi titik' 283I dan 283K di bawah ini.   Namun saya memulai
dengan hasil yang relatif mudah, dan setidaknya sama pentingnya,
menunjukkan (antara lain) bahwa fungsi $f$ yang dapat diintegralkan dapat
(pada dasarnya) dipulihkan dari transformasi Fourier-nya.

\medskip

\noindent}{\bf Lema} Kapan pun $c<d$ di $\Bbb R$,

$$\eqalign{\lim_{a\to\infty}\int_{-a}^{a}
     e^{-iyx}\,\Bover{e^{idy}-e^{icy}}{y}dy
&=2\pi i\text{ jika }c<x<d,\cr
&=\pi i\text{ jika }x=c\text{ atau }x=d,\cr
&=0\text{ jika }x<c\text{ atau } x>d.\cr}$$

\proof{ Kita tahu bahwa untuk setiap $b>0$

\Centerline{$\lim_{a\to\infty}\int_{-a}^{a}\Bover{\sin bx}{x}dx
=\lim_{a\to\infty}\int_{-ab}^{ab}\Bover{\sin t}{t}dt
=\pi$}

\noindent (dengan mengganti $x=t/b$), dan oleh karena itu untuk setiap $b<0$

\Centerline{$\lim_{a\to\infty}\int_{-a}^{a}\Bover{\sin bx}{x}dx
=-\lim_{a\to\infty}\int_{-a}^{a}\Bover{\sin(-b)x}{x}dx
=-\pi$.}

\noindent Sekarang pertimbangkan, untuk $x\in\Bbb R$,

\Centerline{$\lim_{a\to\infty}\int_{-a}^{a}e^{-iyx}
\,\Bover{e^{idy}-e^{icy}}{y}dy$.}

\noindent Pertama, perhatikan bahwa semua integral $\int_{-a}^a$ ada, karena

\Centerline{$\lim_{y\to 0}\Bover{e^{idy}-e^{icy}}{y}=i(d-c)$}

\noindent terbatas, dan integran tersebut kontinu kecuali di
 $0$.   Sekarang kita punya

$$\eqalignno{\int_{-a}^ae^{-iyx}&\Bover{e^{idy}-e^{icy}}{y}dy\cr
&=\int_{-a}^a\Bover{e^{i(d-x)y}-e^{i(c-x)y}}{y}dy\cr
&=\int_{-a}^a\Bover{\cos(d-x)y-\cos(c-x)y}{y}dy
 +i\int_{-a}^a\Bover{\sin(d-x)y-\sin(c-x)y}{y}dy\cr
&=i\int_{-a}^a\Bover{\sin(d-x)y-\sin(c-x)y}{y}dy\cr}$$

\noindent karena $\cos$ adalah fungsi genap, jadi

\Centerline{$\int_{-a}^{a}\Bover{\cos(d-x)y-\cos(c-x)y}{y}dy=0$}

\noindent untuk setiap $a\ge 0$.   (Sekali lagi, integral ini ada
 karena

\Centerline{$\lim_{y\to 0}\Bover{\cos(d-x)y-\cos(c-x)y}{y}=0$.)}

\noindent   Oleh karena itu

$$\eqalign{\lim_{a\to\infty}\int_{-a}^a
e^{-iyx}\Bover{e^{idy}-e^{icy}}{y}dy
&=i\lim_{a\to\infty}\int_{-a}^a\Bover{\sin(d-x)y}{y}dy
   -i\lim_{a\to\infty}\int_{-a}^a\Bover{\sin(c-x)y}{y}dy\cr
&=i\pi-i\pi=0\text{ jika }x<c,\cr
&=i\pi-0=\pi i\text{ jika }x=c,\cr
&=i\pi+i\pi=2\pi i\text{ jika }c<x<d,\cr
&=0+i\pi=\pi i\text{ jika }x=d,\cr
&=-i\pi+i\pi=0\text{ jika }x>d.\cr}$$
}%end of proof of 283E

\leader{283F}{Teorema} Misalkan $f$ adalah fungsi bernilai kompleks yang dapat
diintegralkan pada $\Bbb R$, dan $\varhatf$ transformasi Fourier-nya. Maka
kapan pun $c\le d$ dalam $\Bbb R$,

$$\int_c^df=\Bover{i}{\sqrt{2\pi}}\lim_{a\to\infty}
\int_{-a}^{a}\Bover{e^{icy}-e^{idy}}{y}\varhatf(y)dy.$$

\proof{ Jika $c=d$, hasilnya sepele; kita dapat mengandaikan bahwa $c<d$.

\medskip

{\bf (a)} Dengan menulis

$$\theta_a(x)=\biggerint_{-a}^{a}e^{-iyx}
\,\Bover{e^{idy}-e^{icy}}{y}dy$$

\noindent untuk $x\in\Bbb R$ dan $a\ge 0$, 283E memberitahu kita bahwa

\Centerline{$\lim_{a\to\infty}\theta_a(x)=2\pi i\theta(x)$}

\noindent dengan $\theta=\bover12(\chi[c,d]+\chi\ooint{c,d})$ mengambil
nilai $1$ di dalam interval $[c,d]$, $0$ luar dan nilai
$\bover12$ pada titik akhir.

$$\eqalignno{|\theta_a(x)|
&=|\int_{-a}^a\Bover{\sin(d-x)y-\sin(c-x)y}{y}dy|\cr
\displaycause{lihat pembuktian 283E}
&\le|\int_{-a}^a\Bover{\sin(d-x)y}{y}dy|
  +|\int_{-a}^a\Bover{\sin(c-x)y}{y}dy|
\le 2K\cr}$$

\noindent untuk semua $a\ge 0$ dan $x\in\Bbb R$, dengan $K$ konstanta dari
283Db. Akibatnya $|f\times\theta_a|\le 2K|f|$ di seluruh $\dom f$, untuk
setiap $a\ge 0$, dan (dengan menerapkan Teorema Konvergensi Terdominasi
Lebesgue pada barisan $\sequencen{f\times\theta_{a_n}}$, dengan
$a_n\to\infty$)

\Centerline{$\lim_{a\to\infty}
\int f\times\theta_a=2\pi i\int f\times\theta=2\pi i\int_c^df$.}

\medskip

{\bf (b)} Sekarang pertimbangkan batas dalam pernyataan teorema.
Kita mempunyai

$$\eqalignno{\int_{-a}^a
\bover{e^{icy}-e^{idy}}{y}\varhatf(y)dy
&=\bover1{\sqrt{2\pi}}\int_{-a}^a
\int_{-\infty}^{\infty}
\bover{e^{icy}-e^{idy}}{y}e^{-iyx}f(x)dxdy\cr
&=\bover1{\sqrt{2\pi}}\int_{-\infty}^{\infty}\int_{-a}^a
\bover{e^{icy}-e^{idy}}{y}e^{-iyx}f(x)dydx\cr
&=-\bover1{\sqrt{2\pi}}\int_{-\infty}^{\infty}f\times\theta_a\cr}$$

\noindent oleh teorema Fubini dan Tonelli (252H), dengan menggunakan fakta bahwa
$(e^{icy}-e^{idy})/y$ terbatas untuk melihat bahwa

\Centerline{$\int_{-\infty}^{\infty}\int_{-a}^{a}
  \bigl|\bover{e^{icy}-e^{idy}}{y}e^{-iyx}f(x)\bigr|dydx$}

\noindent berhingga; oleh karena itu, kita dapat menukar urutan integrasi.

$$\eqalign{\bover{i}{\sqrt{2\pi}}\lim_{a\to\infty}\int_{-a}^a
\bover{e^{icy}-e^{idy}}{y}\varhatf(y)dy
&=-\bover{i}{2\pi}\lim_{a\to\infty}
\int_{-\infty}^{\infty}f\times\theta_a\cr
&=-\bover{i}{2\pi}2\pi i\int_c^df
=\int_c^df,\cr}$$

\noindent sesuai kebutuhan.
}%end of proof of 283F

\leader{283G}{Akibat} Jika $f$ dan $g$ adalah fungsi bernilai kompleks yang
dapat diintegralkan pada $\Bbb R$, maka
$\varhatf=\varhat{g}$ jika $f\eae g$.

\proof{ Jika $f\eae g$ maka tentu saja

\Centerline{$\varhatf(y)
=\Bover1{\sqrt{2\pi}}\int_{-\infty}^{\infty}e^{-iyx}f(x)dx
=\Bover1{\sqrt{2\pi}}\int_{-\infty}^{\infty}e^{-iyx}g(x)dx
=\varhat{g}(y)$}

\noindent untuk setiap $y\in\Bbb R$. Sebaliknya, jika
$\varhatf=\varhat{g}$, maka oleh teorema terakhir

\Centerline{$\int_c^df=\int_c^dg$}

\noindent untuk semua $c\le d$, sehingga $f=g$ hampir di mana-mana, menurut 222D.
}%end of proof of 283G

\leader{283H}{Lema} Misalkan $f$ adalah fungsi bernilai kompleks yang dapat
diintegralkan pada $\Bbb R$, dan $\varhatf$ transformasi Fourier-nya. Maka

\Centerline{$\Bover1{\sqrt{2\pi}}\int_{-a}^{a}e^{ixy}\varhatf(y)dy
=\Bover1{\pi}\int_{-\infty}^{\infty}\Bover{\sin a(x-t)}{x-t}f(t)dt
=\Bover1{\pi}\int_{-\infty}^{\infty}\Bover{\sin at}{t}f(x-t)dt$}

\noindent setiap kali $a>0$ dan $x\in\Bbb R$.

\proof{ Kita mempunyai

\Centerline{$\int_{-a}^a\int_{-\infty}^{\infty}|e^{ixy}e^{-iyt}f(t)|dtdy
\le 2a\int_{-\infty}^{\infty}|f(t)|dt<\infty$,}

\noindent jadi (karena fungsi $(t,y)\mapsto e^{ixy}e^{-iyt}f(t)$ adalah
pasti terukur bersama-sama) kita dapat membalikkan urutan integrasi, dan
mendapatkan

$$\eqalign{\Bover1{\sqrt{2\pi}}\int_{-a}^ae^{ixy}\varhatf(y)dy
&=\Bover1{2\pi}\int_{-a}^a\int_{-\infty}^{\infty}
   e^{ixy}e^{-iyt}f(t)dt\,dy\cr
&=\Bover1{2\pi}\int_{-\infty}^{\infty}f(t)\int_{-a}^a
   e^{i(x-t)y}dy\,dt\cr
&=\Bover1{2\pi}\int_{-\infty}^{\infty}
   \Bover{2\sin(x-t)a}{x-t}f(t)dt
=\Bover1{\pi}\int_{-\infty}^{\infty}
   \Bover{\sin au}{u}f(x-u)du,\cr}$$

\noindent menggantikan $t=x-u$.
}%end of proof of 283H

\leader{283I}{Teorema} Misalkan $f$ adalah fungsi bernilai kompleks yang
dapat diintegralkan pada $\Bbb R$, dan anggaplah $f$ terdiferensiasi di
$x\in\Bbb R$. Maka

\Centerline{$f(x)
=\Bover{1}{\sqrt{2\pi}}
\lim_{a\to\infty}\int_{-a}^ae^{ixy}\varhatf(y)dy
=\Bover{1}{\sqrt{2\pi}}
\lim_{a\to\infty}\int_{-a}^ae^{-ixy}\varcheckf(y)dy$.}

\proof{ Tetapkan $g(u)=f(x)$ jika $|u|\le 1$, dan $0$ jika tidak;
perhatikan bahwa
$\lim_{u\to 0}\bover1u(f(x-u)-g(u))=-f'(x)$ terbatas, sehingga ada
suatu $\delta\in\ocint{0,1}$ sedemikian sehingga

\Centerline{$K
=\sup_{0<|u|\le\delta}\bigl|\Bover{f(x-u)-g(u)}{u}\bigr|<\infty$.}

\noindent Akibatnya

$$\eqalign{\int_{-\infty}^{\infty}\bigl|\Bover{f(x-u)-g(u)}{u}\bigr|du
&\le\Bover1{\delta}\int_{-\infty}^{-\delta}|f(x-u)|du
   +\Bover1{\delta}\int_{-1}^1|g|\cr
&{\hskip 8em plus 0pt minus 0pt}
   +\int_{-\delta}^{\delta}K
   +\Bover1{\delta}\int_{\delta}^{\infty}|f(x-u)|du\cr
&\le\Bover{1}{\delta}\int_{-\infty}^{\infty}|f|
   + \Bover2{\delta}|f(x)|+ 2\delta K
<\infty.\cr}$$

\noindent Oleh lema Riemann-Lebesgue (282Fb lagi),

\Centerline{$\lim_{a\to\infty}\int_{-\infty}^{\infty}
\Bover{\sin au}{u}(f(x-u)-g(u))du=0$.}

\noindent Jika kita sekarang memeriksa $\int\bover{\sin au}{u}g(u)du$, kita mendapatkan

$$\eqalign{\int_{-\infty}^{\infty}\Bover{\sin au}{u}g(u)du
&=\int_{-1}^1\Bover{\sin au}{u}f(x)du
=f(x)\int_{-a}^a\Bover{\sin v}{v}dv,\cr}$$

\noindent dengan mengganti $u=v/a$.   Jadi kita peroleh

$$\eqalign{\lim_{a\to\infty}\int_{-\infty}^{\infty}\Bover{\sin
au}{u}f(x-u)du
&=\lim_{a\to\infty}\int_{-\infty}^{\infty}\Bover{\sin au}{u}g(u)du\cr
&=\lim_{a\to\infty}f(x)\int_{-a}^a\Bover{\sin v}{v}dv
=\pi f(x),\cr}$$

\noindent oleh 283Da.   Oleh karena itu

$$\Bover{1}{\sqrt{2\pi}}
\lim_{a\to\infty}\int_{-a}^ae^{ixy}\varhatf(y)dy
=\Bover{1}{\pi}\lim_{a\to\infty}\int_{-\infty}^{\infty}\Bover{\sin
au}{u}f(x-u)du
= f(x),$$

\noindent menggunakan 283H.    Adapun kesetaraan kedua,

$$\eqalignno{\Bover1{\sqrt{2\pi}}\lim_{a\to\infty}\int_{-a}^a
e^{-ixy}\varcheckf(y)dy
&=\Bover1{\sqrt{2\pi}}\lim_{a\to\infty}\int_{-a}^a
e^{-ixy}\varhatf(-y)dy\cr
&=\Bover1{\sqrt{2\pi}}\lim_{a\to\infty}\int_{-a}^a
e^{ixu}\varhatf(u)du
=f(x)\cr}$$

\noindent (memasukkan $y=-u$).
}%end of proof of 283I

\cmmnt{\medskip

\noindent{\bf Keterangan} Bandingkan 282L.}%end of comment

\leader{283J}{Akibat} Misalkan $f:\Bbb R\to\Bbb C$ adalah fungsi yang dapat
diintegralkan sedemikian sehingga $f$ terdiferensiasi dan $\varhatf$ dapat
diintegralkan. Maka $f=(\varhatf)\varspcheck=(\varcheckf)\varsphat$.

\proof{ Karena $\varhatf$ dapat diintegralkan,

\Centerline{$\varhatf\varcheck{\phantom{f}}(x)=\lim_{a\to\infty}
\Bover1{\sqrt{2\pi}}\int_{-a}^ae^{ixy}\varhatf(y)dy = f(x)$}

\noindent untuk setiap $x\in\Bbb R$.  Demikian pula,

\Centerline{$\varcheckf\varhat{\phantom{f}}(x)=\lim_{a\to\infty}
\Bover1{\sqrt{2\pi}}\int_{-a}^ae^{-ixy}\varcheckf(y)dy = f(x).$}
}%end of proof of 283J

\cmmnt{\medskip

\noindent{\bf Catatan} Lihat juga 283Wk di bawah.
}

\leader{283K}{}\cmmnt{ Proposisi berikutnya memberikan kelas fungsi
yang padanya akibat terakhir dapat diterapkan.

\medskip

\noindent}{\bf Proposisi} Misalkan $f$ adalah fungsi dari $\Bbb R$ ke $\Bbb C$
yang terdiferensiasi dua kali dan $f$, $f'$, serta $f''$ semuanya dapat
diintegralkan. Maka $\varhatf$
dapat diintegralkan.

\proof{ Karena $f'$ dan $f''$ dapat diintegralkan, $f$ dan $f'$
kontinu mutlak pada interval terbatas mana pun (225L).   Jadi dengan 283Ci
kita punya

\Centerline{$(f'')\varsphat(y)=iy(f')\varsphat(y)=-y^2\varhatf(y)$}

\noindent untuk setiap $y\in\Bbb R$. Pada saat yang sama, menurut
283Cf-283Cg, $(f'')\varsphat$ dan $\varhatf$ harus terbatas; katakanlah
$|\varhatf(y)|+|(f'')\varsphat(y)|\le K$ untuk setiap $y\in\Bbb R$.   Sekarang

\Centerline{$|\varhatf(y)|\le\Bover{K}{1+y^2}$}

\noindent untuk setiap $y$, sehingga

\Centerline{$\int_{-\infty}^{\infty}|\varhatf|
\le K\int_{-\infty}^{-1}\Bover{1}{y^2}dy
  +2K+K\int_{1}^{\infty}\Bover{1}{y^2}dy
=4K<\infty$.}
}%end of proof of 283K

\cmmnt{\medskip

\noindent{\bf Keterangan} Bandingkan 282Rb.}

\leader{283L}{}\cmmnt{ Sekarang saya beralih ke hasil yang bersesuaian dengan 282O,
menggunakan pendekatan yang sedikit berbeda.

\medskip

\noindent}{\bf Teorema} Misalkan $f$ adalah fungsi bernilai kompleks yang
dapat diintegralkan pada $\Bbb R$, dengan transformasi Fourier $\varhatf$ dan
transformasi Fourier invers $\varcheckf$, dan anggaplah $f$ memiliki variasi
terbatas pada suatu lingkungan dari $x\in\Bbb R$. Tetapkan
$a=\lim_{t\in\dom f,t\uparrow x}f(t)$,
$b=\lim_{t\in\dom f,t\downarrow x}f(t)$. Maka

\Centerline{$\Bover1{\sqrt{2\pi}}\lim_{\gamma\to\infty}
     \int_{-\gamma}^{\gamma}e^{ixy}\varhatf(y)dy
=\Bover1{\sqrt{2\pi}}\lim_{\gamma\to\infty}
     \int_{-\gamma}^{\gamma}e^{-ixy}\varcheckf(y)dy
=\Bover12(a+b)$.}

\proof{{\bf (a)} Limit $\lim_{t\in\dom f,t\uparrow x}f(t)$ dan
$\lim_{t\in\dom f,t\downarrow x}f(t)$ ada karena $f$ memiliki variasi
terbatas di dekat $x$ (224F). Ingat dari 283Db bahwa terdapat konstanta
$K<\infty$ sedemikian sehingga

\Centerline{$|\int_{\gamma}^{\delta}\Bover{\sin cx}{x}dx|\le K$}

\noindent setiap kali $\gamma\le \delta$ dan $c\in\Bbb R$.

\medskip

{\bf (b)} Misalkan $\epsilon>0$.   Hipotesisnya adalah ada beberapa
$\delta>0$ sehingga $\Var_{[x-\delta,x+\delta]}(f)<\infty$.
Akibatnya

\Centerline{$\lim_{\eta\downarrow 0}\Var_{\ocint{x,x+\eta}}(f)
=\lim_{\eta\downarrow 0}\Var_{\coint{x-\eta,x}}(f)=0$}

\noindent (224E).   Oleh karena itu ada $\eta>0$ sedemikian sehingga

\Centerline{$\max(\Var_{\coint{x-\eta,x}}(f),
               \Var_{\ocint{x,x+\eta}}(f))\le\epsilon$.}

\noindent Tentu saja

\Centerline{$|f(t)-f(u)|\le\Var_{\coint{x-\eta,x}}(f)\le\epsilon$}

\noindent kapan pun $t$, $u\in\dom f$ dan $x-\eta\le t\le u<x$, sehingga kita
akan memiliki

\Centerline{$|f(t)-a|\le\epsilon$ untuk setiap
$t\in\dom f\cap\coint{x-\eta,x}$,}

\noindent dan demikian pula

\Centerline{$|f(t)-b|\le\epsilon$ setiap kali $t\in\dom f
\cap\ocint{x,x+\eta}$.}

\medskip

{\bf (c)} Sekarang tetapkan

\Centerline{$g_1(t)=f(t)$ ketika $t\in\dom f$ dan $|x-t|>\eta$, $0$
sebaliknya,}

\Centerline{$g_2(t)=a$ ketika $x-\eta\le t<x$, $b$ ketika $x<t\le x+\eta$,
$0$ sebaliknya,}

\Centerline{$g_3=f-g_1-g_2$.}

\noindent Maka $f=g_1+g_2+g_3$;  setiap $g_j$ dapat diintegralkan;  $g_1$ adalah
nol di lingkungan $x$;

\Centerline{$\sup_{t\in\dom g_3,t\ne x}|g_3(t)|\le\epsilon$,}

\Centerline{$\Var_{\coint{x-\eta,x}}(g_3)\le\epsilon$,
\quad $\Var_{\ocint{x,x+\eta}}(g_3)\le\epsilon$.}

\medskip

{\bf (d)} Pertimbangkan tiga bagian $g_1$, $g_2$, $g_3$ secara terpisah.

\medskip

\quad{\bf (i)} Untuk yang pertama, kita punya

\Centerline{$\lim_{\gamma\to\infty}\Bover1{\sqrt{2\pi}}
   \int_{-\gamma}^{\gamma}e^{ixy}\varhat{g}_1(y)dy=0$}

\noindent oleh 283I.

\medskip

\quad{\bf (ii)} Selanjutnya,

$$\eqalignno{\Bover{1}{\sqrt{2\pi}}
        \int_{-\gamma}^{\gamma}e^{ixy}\varhat{g}_2(y)dy
&=\Bover1{\pi}\int_{-\infty}^{\infty}
        \Bover{\sin(x-t)\gamma}{x-t}g_2(t)dt\cr
\noalign{\noindent (menurut 283H)}
&=\Bover{a}{\pi}\int_{x-\eta}^{x}\Bover{\sin(x-t)\gamma}{x-t}dt
   +\Bover{b}{\pi}\int_{x}^{x+\eta}\Bover{\sin(x-t)\gamma}{x-t}dt\cr
&=\Bover{a}{\pi}\int_0^{\gamma\eta}\Bover{\sin u}{u}du
   +\Bover{b}{\pi}\int_{0}^{\gamma\eta}\Bover{\sin u}{u}du\cr
\noalign{\noindent (dengan substitusi $t=x-\bover{1}{\gamma}u$ pada integral pertama,
$t=x+\bover1{\gamma}u$ pada integral kedua)}
&\to \Bover{a+b}{2}\text{ ketika }\gamma\to\infty\cr}$$

\noindent oleh 283Da.

\medskip

\quad{\bf (iii)} Untuk yang ketiga, kita mempunyai, untuk setiap $\gamma>0$,

$$\eqalignno{\bigl|\Bover1{\sqrt{2\pi}}\int_{-\gamma}^{\gamma}e^{ixy}
     \varhat{g}_3(y)dy\bigr|
&=\Bover1{\pi}\bigl|\int_{-\infty}^{\infty}
     \Bover{\sin(x-t)\gamma}{x-t}g_3(t)dt\bigr|
=\Bover1{\pi}\bigl|\int_{-\infty}^{\infty}
     \Bover{\sin t\gamma}{t}g_3(x-t)dt\bigr|\cr
&\le\Bover1{\pi}\bigl|\int_{-\eta}^0
     \Bover{\sin t\gamma}{t}g_3(x-t)dt\bigr|
   +\Bover1{\pi}\bigl|\int_0^{\eta}
     \Bover{\sin t\gamma}{t}g_3(x-t)dt\bigr|\cr
&\le \Bover{K}{\pi}\Bigl(\sup_{t\in\dom g_3\cap\ooint{x-\eta,x}}|g_3(t)|
   +\Var_{\ooint{x-\eta,x}}(g_3)\cr
  &{\hskip 8em plus 0pt minus 0pt}
   +\sup_{t\in\dom g_3\cap\ooint{x,x+\eta}}|g_3(t)|
   +\Var_{\ooint{x,x+\eta}}(g_3)\Bigr)\cr
&\le 4\epsilon\Bover{K}{\pi},\cr}$$

\noindent dengan menggunakan 224J lagi untuk memperkirakan integral
berdasarkan variasi dan supremum $g_3$ serta integral
$\bover{\sin\gamma t}{t}$ pada subinterval.

\medskip

{\bf (e)} Oleh karena itu, kita mempunyai

$$\eqalign{\limsup_{\gamma\to\infty}\bigl|\Bover1{\sqrt{2\pi}}
        \int_{-\gamma}^{\gamma}
   &e^{ixy}\varhatf(y)dy-\Bover{a+b}{2}\bigr|\cr
&\le\limsup_{\gamma\to\infty}\Bover1{\sqrt{2\pi}}
     \bigl|\int_{-\gamma}^{\gamma}e^{ixy}\varhat{g}_1(y)dy\bigr|\cr
&{\hskip 6em plus 0pt minus 0pt}
     +\limsup_{\gamma\to\infty}\bigl|\Bover1{\sqrt{2\pi}}
     \int_{-\gamma}^{\gamma}e^{ixy}\varhat{g}_2(y)dy
     -\Bover{a+b}{2}\bigr|\cr
&{\hskip 6em plus 0pt minus 0pt}
     +\limsup_{\gamma\to\infty}\Bover1{\sqrt{2\pi}}
     \bigl|\int_{-\gamma}^{\gamma}e^{ixy}\varhat{g}_3(y)\,dy\bigr|\cr
&\le 0+0+\Bover{4K}{\pi}\epsilon\cr}$$

\noindent dengan perhitungan pada (d).   Karena $\epsilon$ bersifat arbitrer,

\Centerline{$\lim_{\gamma\to\infty}\Bover1{\sqrt{2\pi}}
     \int_{-\gamma}^{\gamma}e^{ixy}\varhatf(y)dy-\Bover{a+b}{2}=0$.}

\medskip

{\bf (f)} Ini adalah paruh pertama teorema.   Namun tentu saja bagian kedua
segera menyusul, karena

$$\eqalign{\Bover1{\sqrt{2\pi}}\lim_{\gamma\to\infty}
    \int_{-\gamma}^{\gamma}e^{-ixy}\varcheckf(y)dy
&=\Bover1{\sqrt{2\pi}}\lim_{\gamma\to\infty}
    \int_{-\gamma}^{\gamma}e^{-ixy}\varhatf(-y)dy\cr
&=\Bover1{\sqrt{2\pi}}\lim_{\gamma\to\infty}
    \int_{-\gamma}^{\gamma}e^{ixy}\varhatf(y)dy
=\Bover{a+b}2.\cr}$$
}%end of proof of 283L

\cmmnt{\medskip

\noindent{\bf Catatan} Anda akan melihat bahwa argumen ini menggunakan beberapa ide
yang sama seperti yang ada di 282O--282P.   Argumen ini lebih langsung karena
(i) saya tidak menggunakan konsep apa pun yang bersesuaian dengan jumlah Fej\'er
(meskipun ada konsep yang sangat sesuai untuk tujuan ini; lihat 283Xf);
(ii) saya tidak repot-repot memberikan hasil tentang konvergensi seragam
integral Fej\'er ketika $f$ kontinu dan mempunyai variasi terbatas (283Xj);
dan (iii) saya tidak menyinggung pentingnya fakta bahwa jika $f$ mempunyai
variasi terbatas, maka $\sup_{y\in\Bbb R}|y\varhatf(y)|<\infty$ (283Xk).
}%end of comment

\leader{283M}{}\cmmnt{ Sesuai dengan 282Q, kami memiliki yang berikut ini.

\medskip

\noindent}{\bf Teorema} Misalkan $f$ dan $g$ adalah fungsi bernilai kompleks
yang dapat diintegralkan pada $\Bbb R$, dan $f*g$ produk konvolusinya,
ditentukan dengan menetapkan

\Centerline{$(f*g)(x)=\int_{-\infty}^{\infty}f(t)g(x-t)dt$}

\noindent setiap kali ini didefinisikan\cmmnt{ (255E)}.   Maka

\Centerline{$(f*g)\varsphat(y)
=\sqrt{2\pi}\varhatf(y)\varhat{g}(y)$,\quad
$(f*g)\varspcheck(y)
=\sqrt{2\pi}\varcheckf(y)\varcheck{g}(y)$}

\noindent untuk setiap $y\in\Bbb R$.

\proof{ Untuk setiap $y$,

$$\eqalignno{(f*g)\varsphat(y)
&=\Bover1{\sqrt{2\pi}}\int_{-\infty}^{\infty}e^{-iyx}(f*g)(x)dx\cr
&=\Bover1{\sqrt{2\pi}}\int_{-\infty}^{\infty}\int_{-\infty}^{\infty}
e^{-iy(t+u)}f(t)g(u)dtdu\cr
\noalign{\noindent (dengan menggunakan 255G)}
&=\Bover1{\sqrt{2\pi}}\int_{-\infty}^{\infty}e^{-iyt}f(t)dt
\int_{-\infty}^{\infty}e^{-iyu}g(u)du
=\sqrt{2\pi}\varhatf(y)\varhat{g}(y).\cr}$$

\noindent Sekarang, tentu saja,

\Centerline{$(f*g)\varspcheck(y)
=(f*g)\varsphat(-y)
=\sqrt{2\pi}\varhatf(-y)\varhat{g}(-y)
=\sqrt{2\pi}\varcheckf(y)\varcheck{g}(y)$.}
}%end of proof of 283M

\leader{283N}{}\cmmnt{ Saya menunjukkan cara menghitung transformasi Fourier
khusus, yang mana
akan digunakan berulang kali di bagian berikutnya.

\medskip

\noindent}{\bf Lema} Untuk $\sigma>0$, atur
$\psi_{\sigma}(x)=\bover1{\sigma\sqrt{2\pi}}e^{-x^2/2\sigma^2}$ untuk
$x\in\Bbb R$.   Maka transformasi Fourier dan transformasi Fourier invers
adalah

\Centerline{$\varhat{\psi}_{\sigma}=\varcheck{\psi}_{\sigma}
=\Bover1{\sigma}\psi_{1/\sigma}$.}

\noindent Khususnya, $\varhat{\psi}_1=\psi_1$.

\proof{{\bf (a)} Saya memulai dengan kasus khusus $\sigma=1$, menggunakan deret
Maclaurin

\Centerline{$e^{-iyx}=\sum_{k=0}^{\infty}\Bover{(-iyx)^k}{k!}$}

\noindent dan ekspresi untuk
$\int_{-\infty}^{\infty}x^ke^{-x^2/2}dx$ dari \S263.

Tetapkan $y\in\Bbb R$.    Menulis

\Centerline{$g_k(x)=\Bover{(-iyx)^k}{k!}e^{-x^2/2}$,\quad
$h_n(x)=\sum_{k=0}^ng_k(x)$, \quad$h(x)=e^{|yx|-x^2/2}$,}

\noindent kita melihat bahwa

\Centerline{$|g_k(x)|\le \Bover{|yx|^k}{k!}e^{-x^2/2}$,}

\noindent sehingga

\Centerline{$|h_n(x)|\le\sum_{k=0}^{\infty}|g_k(x)|
\le e^{|yx|}e^{-x^2/2}=h(x)$}

\noindent untuk setiap $n$;  terlebih lagi, $h$ dapat diintegralkan, karena
$|h(x)|\le e^{-|x|}$ setiap kali $|x|\ge 2(1+|y|)$.   Akibatnya, menggunakan
Teorema Konvergensi Dominasi Lebesgue,

$$\eqalignno{\varhat{\psi}_1(y)
&=\Bover1{{2\pi}}\int_{-\infty}^{\infty}
  \lim_{n\to\infty}h_n\
=\Bover1{{2\pi}}\lim_{n\to\infty}\int_{-\infty}^{\infty}
  h_n\cr
&=\Bover1{{2\pi}}\sum_{k=0}^{\infty}
  \int_{-\infty}^{\infty}g_k
=\Bover1{{2\pi}}\sum_{k=0}^{\infty}\Bover{(-iy)^k}{k!}
  \int_{-\infty}^{\infty}x^ke^{-x^2/2}dx\cr
&=\Bover1{{2\pi}}\sum_{j=0}^{\infty}\Bover{(-iy)^{2j}}{(2j)!}
  \Bover{(2j)!}{2^jj!}\sqrt{2\pi}\cr
\noalign{\noindent (menurut 263H)}
&=\Bover1{\sqrt{2\pi}}\sum_{j=0}^{\infty}\Bover{(-y^2)^j}{2^jj!}
=\Bover1{\sqrt{2\pi}}e^{-y^2/2}
=\psi_1(y),\cr}$$

\noindent seperti yang diklaim.

\medskip

{\bf (b)} Untuk kasus umum
$\psi_{\sigma}(x)=\Bover1{\sigma}\psi_1(\Bover{x}{\sigma})$, sehingga

\Centerline{$\varhat{\psi}_{\sigma}(y)
=\Bover1{\sigma}\cdot\sigma\varhat{\psi}_1(\sigma y)
=\Bover1\sigma\psi_{1/\sigma}(y)$}

\noindent oleh 283Ce.   Tentu saja sekarang kita punya

\Centerline{$\varcheck{\psi}_{\sigma}(y)=\varhat{\psi}_{\sigma}(-y)
=\Bover1{\sigma}\psi_{1/\sigma}(y)$}

\noindent karena $\psi_{1/\sigma}$ merupakan fungsi genap.
}%end of proof of 283N


\leader{283O}{}\cmmnt{ Untuk mengarahkan ide pada bagian selanjutnya, saya
memberikan fakta yang sangat sederhana berikut.

\medskip

\noindent}{\bf Proposisi} Misalkan $f$ dan $g$ adalah dua fungsi bernilai
kompleks yang dapat diintegralkan pada $\Bbb R$. Maka
$\int_{-\infty}^{\infty}f\times\varhat{g}
=\int_{-\infty}^{\infty}\varhatf\times g$ dan
$\int_{-\infty}^{\infty}f\times\varcheck{g}
  =\int_{-\infty}^{\infty}\varcheckf\times g$.

\proof{  Tentu saja

\Centerline{$\int_{-\infty}^{\infty}\int_{-\infty}^{\infty}
|e^{-ixy}f(x)g(y)|dxdy
=\int_{-\infty}^{\infty}|f|\int_{-\infty}^{\infty}|g|<\infty$,}

\noindent jadi

$$\eqalign{\int_{-\infty}^{\infty}f\times\varhat{g}
&=\Bover1{\sqrt{2\pi}}\int_{-\infty}^{\infty}\int_{-\infty}^{\infty}
f(y)e^{-iyx}g(x)dxdy\cr
&=\Bover1{\sqrt{2\pi}}\int_{-\infty}^{\infty}\int_{-\infty}^{\infty}
f(y)e^{-ixy}g(x)dydx
=\int_{-\infty}^{\infty}\varhatf\times g.\cr}$$

\noindent Untuk separuh proposisi lainnya, ganti setiap
$e^{-ixy}$ dalam argumen dengan $e^{ixy}$.
}%end of proof of 283O

\exercises{\leader{283W}{Dimensi yang lebih tinggi} Saya menawarkan sebarisan latihan
yang dirancang untuk memberikan petunjuk tentang bagaimana pekerjaan bagian ini dapat dilakukan
dalam kasus dimensi
$r$, di mana $r\ge 1$.

\spheader 283Wa Misalkan $f$ adalah fungsi bernilai kompleks yang dapat
diintegralkan dan terdefinisi hampir di mana-mana pada $\BbbR^r$.
{\bf Transformasi Fourier}-nya
adalah fungsi $\varhatf:\BbbR^r\to\Bbb C$ yang ditentukan oleh rumus

\Centerline{$\varhatf(y)
=\Bover1{(\sqrt{2\pi})^r}\int e^{-iy\dotproduct x}f(x)dx$,}

\noindent dengan menulis $y\dotproduct x=\eta_1\xi_1+\ldots+\eta_r\xi_r$ untuk
$x=(\xi_1,\ldots,\xi_r)$ dan $y=(\eta_1,\ldots,\eta_r)\in\BbbR^r$, dan
$\int\ldots dx$ untuk integrasi sehubungan dengan ukuran Lebesgue di
$\BbbR^r$.   Demikian pula, transformasi Fourier {\bf invers} dari $f$ adalah
fungsi $\varcheckf$ yang diberikan oleh

\Centerline{$\varcheckf(y)
=\Bover1{(\sqrt{2\pi})^r}\int e^{iy\dotproduct x}f(x)dx=\varhatf(-y)$.}

\noindent Tunjukkan bahwa, untuk fungsi bernilai kompleks apa pun yang dapat diintegralkan $f$ pada
$\BbbR^r$,

\quad (i) $\varhatf:\BbbR^r\to\Bbb C$ kontinu;

\quad (ii) $\lim_{\|y\|\to\infty}\varhatf(y)=0$, menulis
$\|y\|=\sqrt{y\dotproduct y}$ seperti biasa;

\quad (iii) jika $\int\|x\||f(x)|dx<\infty$, maka $\varhatf$ adalah
terdiferensiasi, dan

\Centerline{$\Pd{}{\eta_j}\varhatf(y)
=-\Bover{i}{(\sqrt{2\pi})^r}\int e^{-iy\dotproduct x}\xi_jf(x)dx$}

\noindent untuk $j\le r$, $y\in\BbbR^r$, dengan selalu menganggap $\xi_j$
sebagai koordinat ke-$j$ dari $x\in\BbbR^r$;

\quad (iv) jika $j\le r$ dan $\pd{f}{\xi_j}$ didefinisikan di mana-mana dan
dapat diintegralkan, maka
$(\pd{f}{\xi_j})\varsphat(y)=i\eta_j\varhatf(y)$
untuk setiap $y\in\Bbb R^r$.   \prooflet{(Gunakan 225L untuk menunjukkan bahwa jika
$e\in\BbbR^r$ adalah vektor satuan, maka
$\gamma\mapsto f(x+\gamma e)$ kontinu mutlak pada setiap interval
yang dibatasi untuk hampir setiap $x$.)}
%283C

\spheader 283Wb Tunjukkan bahwa jika $f_1,\ldots,f_r$ adalah fungsi bernilai kompleks
yang dapat diintegralkan pada $\Bbb R$ dengan transformasi Fourier
$g_1,\ldots,g_r$, dan kita menulis $f(x)=f_1(\xi_1)\ldots f_r(\xi_r)$ untuk
$x=(\xi_1,\ldots,\xi_r)\in\BbbR^r$, maka transformasi Fourier dari $f$
adalah $y\mapsto g_1(\eta_1)\ldots g_r(\eta_r)$.

\spheader 283Wc Misalkan $f$ adalah fungsi bernilai kompleks yang dapat diintegralkan
pada $\BbbR^r$, dan $\varhatf$ merupakan transformasi Fourier-nya.   Jika $c\le d$ di
$\BbbR^r$, tunjukkan bahwa

$$\int_{[c,d]}f
=(\bover{i}{\sqrt{2\pi}})^r
\lim_{\alpha_1,\ldots,\alpha_r\to\infty}
\int_{[-a,a]}\prod_{j=1}^r
\bover{e^{i\gamma_j\eta_j}-e^{i\delta_j\eta_j}}{\eta_j}\varhatf(y)dy,$$

\noindent dengan notasi $a=(\alpha_1,\ldots)$, $c=(\gamma_1,\ldots)$,
$d=(\delta_1,\ldots)$.
%283F

\spheader 283Wd
Misalkan $f$ adalah fungsi bernilai kompleks yang dapat diintegralkan pada $\BbbR^r$, dan
$\varhatf$ transformasi Fourier-nya.   Tunjukkan bahwa jika kita menulis

\Centerline{$B_{\infty}(\tbf{0},a)=\{y:|\eta_j|\le a
  \text{ untuk setiap }j\le r\}$,}

\noindent lalu

\Centerline{$\Bover1{(\sqrt{2\pi})^r}\int_{B_{\infty}(\tbf{0},a)}
e^{ix\dotproduct y}\varhatf(y)dy
=\int \phi_a(t)f(x-t)dt$}

\noindent untuk setiap $a\ge 0$, di mana

\Centerline{$\phi_a(t)
=\Bover1{\pi^r}\prod_{j=1}^r\Bover{\sin a\tau_j}{\tau_j}$}

\noindent untuk $t=(\tau_1,\ldots,\tau_r)\in\BbbR^r$.
%283H

\spheader 283We Tunjukkan bahwa
$\biggerint_{\BbbR^r}\Bover{1}{1+\|x\|^{r+1}}dx<\infty$.

\spheader 283Wf Misalkan $f:\BbbR^r\to\Bbb C$ adalah fungsi yang dapat diintegralkan.
Tunjukkan bahwa jika semua turunan parsial
$\Bover{\partial^kf}{\partial\xi_j^k}$, untuk $k\le r+1$ dan
$j\le r$, didefinisikan hampir di mana-mana dan dapat diintegralkan,
maka $\varhat{f}$ dapat diintegralkan.
%283Wa(iii) 283We 283K

\spheader 283Wg\dvAformerly{2{}83Wj}
Tunjukkan bahwa jika $f$ dan $g$ merupakan fungsi bernilai kompleks
yang dapat diintegralkan pada $\BbbR^r$, maka (mendefinisikan konvolusi seperti pada
255L) $(f*g)\varsphat=(\sqrt{2\pi})^r\varhatf\times\varhat{g}$.
%283M

\spheader 283Wh Misalkan $f$ dan $g$ adalah fungsi bernilai kompleks yang dapat
diintegralkan pada $\BbbR^r$. Tunjukkan bahwa
$f*\varcheck{g}=(\sqrt{2\pi})^r(\varhatf\times g)\varspcheck$.
%283M 283Wg

\spheader 283Wi\dvAformerly{2{}83We} Untuk $\sigma>0$, tentukan
$\psi_{\sigma}:\Bbb R^r\to\Bbb C$ dengan mengatur

\Centerline{$\psi_{\sigma}(x)=\Bover1{(\sigma\sqrt{2\pi})^r}
e^{-x\dotproduct x/2\sigma^2}$,
\quad$(\varhat{\psi}_{\sigma})\varspcheck=\psi_{\sigma}$.}

\noindent untuk setiap $x\in\BbbR^r$.   Tunjukkan bahwa

\Centerline{$\varhat{\psi}_{\sigma}
=\varcheck{\psi}_{\sigma}=\Bover1{\sigma^r}\psi_{1/\sigma}$.}
%283N

\spheader 283Wj Dengan mendefinisikan $\psi_{\sigma}$ seperti pada bagian (i) di atas, tunjukkan bahwa
$\lim_{\sigma\to 0}(f*\psi_{\sigma})(x)=f(x)$ setiap kali $x\in\BbbR^r$ dan
$f:\BbbR^r\to\Bbb C$ kontinu dan dapat diintegralkan atau terbatas.
(Lih.\ 261Ye, 262Yi.)
%283Wi

\spheader 283Wk Tunjukkan jika $f:\BbbR^r\to\Bbb C$ kontinu
dan dapat diintegralkan, dan $\varhatf$ juga dapat diintegralkan
$f=\varhatf\varcheck{\phantom{f}}$.
({\it Petunjuk\/}: Tunjukkan bahwa keduanya sama di setiap titik
$\lim_{\sigma\to 0}(\sqrt{2\pi})^r
  (\varhatf\times\varhat{\psi}_{\sigma})\varspcheck$.)
%283Wj 283Wh 283J

\spheader 283Wl Tunjukkan bahwa jika $f$ dan $g$ merupakan fungsi bernilai kompleks
yang dapat diintegralkan pada $\BbbR^r$, maka $\int
f\times\varhat{g}=\int\varhatf\times g$.

\spheader 283Wm(i) Tunjukkan bahwa
$\int_{2k\pi}^{2(k+1)\pi}\bover{\sin t}{t\sqrt{t}}dt>0$ untuk setiap
$k\in\Bbb N$, dan karenanya
$\int_0^{\infty}\bover{\sin t}{t\sqrt{t}}dt>0$.

\quad{(ii)} Tetapkan $f_1(\xi)=1/\sqrt{|\xi|}$ untuk $0<|\xi|\le 1$, dan $0$ untuk
 $\xi$ lainnya.   Tunjukkan bahwa
$\lim_{a\to\infty}\bover1{\sqrt{a}}\int_{-a}^a\varhatf_1(\eta)d\eta$
ada di $\Bbb R$ dan lebih besar dari $0$.

\quad{(iii)} Buatlah fungsi $f_2$ yang dapat diintegralkan, nol pada suatu lingkungan
di $0$, sehingga terdapat tak hingga banyak $m\in\Bbb N$
yang $|\int_{-m}^m\varhatf_2(\eta)d\eta|\ge\bover1{\sqrt{m}}$.
\Hint{ambil $f_2(\xi)=2^{-k}\sin m_k\xi$ untuk $k+1\le\xi<k+2$, untuk
barisan $\sequence{k}{m_k}$ yang meningkat cukup cepat.}

\quad{(iv)} Tetapkan $f(x)=f_1(\xi_1)f_2(\xi_2)$ untuk $x\in\BbbR^2$.
Tunjukkan bahwa $f$ dapat diintegralkan, bahwa $f$ bernilai nol pada suatu
lingkungan dari $\tbf{0}$, tetapi bahwa

\Centerline{$\limsup_{a\to\infty}\Bover1{2\pi}|\int_{B_{\infty}(\tbf{0},
a)}\varhatf(y)dy|>0$,}

\noindent mendefinisikan $B_{\infty}$ seperti di 283Wd.
}%end of 283W

\exercises{
\leader{283X}{Latihan dasar (a)}
%\spheader 283Xa
Konfirmasikan bahwa enam alternatif definisi
dari transformasi $\varhatf$, $\varcheckf$ yang ditawarkan di 283B
semuanya mengarah pada teori yang sama;  temukan konstanta yang terlibat dalam versi
baru dari 283Ch, 283Ci, 283L, 283M dan 283N.
%283B

\spheader 283Xb Jika kita mendefinisikan ulang $\varhatf(y)$ menjadi
$\alpha\int_{-\infty}^{\infty}e^{i\beta xy}f(x)dx$,
$\varcheckf(y)$ jadi apa?
%283B

\spheader 283Xc Tunjukkan bahwa hampir semua faktor $2\pi$ akan hilang
dari teorema-teorema bagian ini jika kita mendefinisikan ukuran $\nu$ pada
$\Bbb R$ dengan menetapkan $\nu E=\bover1{\sqrt{2\pi}}\mu E$ untuk setiap
himpunan Lebesgue terukur $E$, dengan $\mu$ ukuran Lebesgue, dan menulis

\Centerline{$\varhatf(y)=\int_{-\infty}^{\infty}e^{-iyx}f(x)\nu(dx)$,
\quad$\varcheckf(y)=\int_{-\infty}^{\infty}e^{iyx}f(x)\nu(dx)$,}

\Centerline{$(f*g)(x)=\int_{-\infty}^{\infty}f(t)g(x-t)\nu(dt)$.}

\noindent Apa itu
$\lim_{a\to\infty}\int_{-a}^a\bover{\sin t}{t}\nu(dt)$?
%283B+

\sqheader 283Xd Misalkan $f$ adalah fungsi bernilai kompleks yang dapat diintegralkan
pada $\Bbb R$, dengan transformasi Fourier $\varhatf$.   Tunjukkan bahwa
(i) jika $g(x)=f(-x)$ kapan pun ini
didefinisikan, lalu $\varhat{g}(y)=\varhatf(-y)$ untuk
setiap $y\in\Bbb R$;
(ii) jika $g(x)=\overline{f(x)}$ kapan pun ini didefinisikan, maka
$\varhat{g}(y)=\overline{\varhatf(-y)}$ untuk setiap $y$.
%283C

\spheader 283Xe Misalkan $f$ adalah fungsi bernilai kompleks yang dapat diintegralkan
pada $\Bbb R$, dengan transformasi Fourier $\varhatf$.   Tunjukkan bahwa

\Centerline{$\int_c^d\varhatf(y)dy=\Bover{i}{\sqrt{2\pi}}
\int_{-\infty}^{\infty}\Bover{e^{-idx}-e^{-icx}}{x}f(x)dx$}

\noindent setiap kali $c\le d$ di $\Bbb R$.
%283F

\sqheader 283Xf Untuk fungsi bernilai kompleks yang dapat diintegralkan $f$ pada
$\Bbb R$, misalkan integral {\bf Fej\'er} didefinisikan sebagai

\Centerline{$\sigma_c(x)
=\Bover1{c\sqrt{2\pi}}\int_0^c
  \bigl(\int_{-a}^{a}e^{ixy}\varhatf(y)dy\bigr)da$}

\noindent untuk $c>0$.   Tunjukkan bahwa

\Centerline{$\sigma_c(x)=
\Bover1{\pi}\int_{-\infty}^{\infty}
  \Bover{1-\cos ct}{ct^2}f(x-t)dt$.}
%283F

\spheader 283Xg Tunjukkan bahwa
$\biggerint_{-\infty}^{\infty}\Bover{1-\cos at}{at^2}dt=\pi$ untuk setiap
$a>0$. \Hint{Integrasikan secara parsial dan gunakan 283Da.} Tunjukkan bahwa

\Centerline{$\lim_{a\to\infty}\int_{\delta}^{\infty}
         \Bover{1-\cos at}{at^2}dt
=\lim_{a\to\infty}\sup_{t\ge\delta}\Bover{1-\cos at}{at^2}
=0$}

\noindent untuk setiap $\delta>0$.
%283Xf, 283D, 283F

\wheader{283Xh}{0}{0}{0}{72pt}

\spheader 283Xh Misalkan $f$ adalah fungsi bernilai kompleks yang dapat diintegralkan
pada $\Bbb R$, dan tentukan integral Fej\'er $\sigma_a$ seperti pada 283Xf
di atas.   Tunjukkan bahwa jika $x\in\Bbb R$, $c\in\Bbb C$ sedemikian sehingga

\Centerline{$\lim_{\delta\downarrow 0}\Bover1{\delta}\int_0^{\delta}
    |f(x+t)+f(x-t)-2c|dt=0$,}

\noindent lalu $\lim_{a\to\infty}\sigma_a(x)=c$.   \Hint{mengadaptasi
argumen 282H.}
%283Xf, 283Xg

\sqheader 283Xi Misalkan $f$ adalah fungsi bernilai kompleks yang dapat diintegralkan
pada $\Bbb R$, dan tentukan integral Fej\'er $\sigma_a$ seperti pada 283Xf
di atas.   Tunjukkan bahwa $f(x)=\lim_{a\to\infty}\sigma_a(x)$ untuk hampir setiap
$x\in\Bbb R$.
%283Xf, 283Xh

\spheader 283Xj Misalkan $f:\Bbb R\to\Bbb C$ adalah fungsi bernilai kompleks
yang kontinu, dapat diintegralkan, dan memiliki variasi terbatas; definisikan
integral Fej\'er-nya $\sigma_a$ seperti pada 283Xf di atas. Tunjukkan bahwa
$f(x)=\lim_{a\to\infty}\sigma_a(x)$ seragam untuk $x\in\Bbb R$.
%283Xf, 283Xi

\sqheader 283Xk Misalkan $f$ adalah fungsi bernilai kompleks yang dapat diintegralkan dan
mempunyai variasi terbatas pada $\Bbb R$, dan $\varhatf$ transformasi Fourier-nya.
Tunjukkan bahwa $\sup_{y\in\Bbb R}|y\varhatf(y)|<\infty$.
%283F

\spheader 283Xl Misalkan $f$ dan $g$ adalah fungsi
bernilai kompleks yang dapat diintegralkan pada $\Bbb R$.   Tunjukkan bahwa
$f*\varcheck{g}=\sqrt{2\pi}(\varhatf\times g)\varspcheck$.
%283M

\spheader 283Xm Misalkan $f$ adalah fungsi bernilai kompleks yang dapat diintegralkan
pada $\Bbb R$, dan tetapkan $x\in\Bbb R$.    Tetapkan

\Centerline{$\hat f_x(y)=\int_{-\infty}^{\infty}f(t)\cos y(x-t)dt$}

\noindent untuk $y\in\Bbb R$.   Tunjukkan bahwa

\quad (i) jika $f$ terdiferensiasi di $x$,

\Centerline{$f(x)
=\Bover1{\pi}\lim_{a\to\infty}\int_{0}^a\hat f_x(y)dy$;}

\quad (ii) jika terdapat lingkungan $x$ di mana $f$ mempunyai variasi terbatas, maka

\Centerline{$\Bover1{\pi}\lim_{a\to\infty}\int_{0}^a\hat f_x(y)dy
=\Bover12(\lim_{t\in\dom f,t\uparrow x}f(t)
 +\lim_{t\in\dom f,t\downarrow x}f(t))$;}

\quad (iii) jika $f$ terdiferensiasi dua kali dan $f'$, $f''$ dapat diintegralkan
maka $\hat f_x$ dapat diintegralkan dan
$f(x)=\bover1{\pi}\int_0^{\infty}\hat f_x$.
(Rumusnya

\Centerline{$f(x)=\Bover1{\pi}\int_{0}^{\infty}\bigl(
\int_{-\infty}^{\infty}f(t)\cos y(x-t)dt\bigr)dy$,}

\noindent berlaku untuk fungsi $f$ tersebut; rumus ini disebut rumus integral
{\bf Fourier}.)

\spheader 283Xn Tunjukkan bahwa jika $f$ adalah fungsi bernilai kompleks
dengan variasi terbatas, terdefinisi hampir di mana-mana pada $\Bbb R$, dan
konvergen ke $0$ (sepanjang domainnya) pada $\pm\infty$, maka

\Centerline{$g(y)
=\Bover1{\sqrt{2\pi}}\lim_{a\to\infty}\int_{-a}^ae^{-iyx}f(x)dx$}

\noindent terdefinisi dalam $\Bbb C$ untuk setiap $y\ne 0$, dan limitnya
seragam pada setiap daerah yang berjarak positif dari $0$.

\spheader 283Xo Misalkan $f$ adalah fungsi bernilai kompleks
yang dapat diintegralkan pada $\Bbb R$.   Tetapkan

\Centerline{$\varhatf_c(y)=\Bover1{\sqrt{2\pi}}\int_{-\infty}^{\infty}
\cos yx\,f(x)dx$,\quad
$\varhatf_s(y)=\Bover1{\sqrt{2\pi}}\int_{-\infty}^{\infty}
\sin yx\,f(x)dx$}

\noindent untuk $y\in\Bbb R$.   Tunjukkan bahwa

\Centerline{$\Bover1{\sqrt{2\pi}}\int_{-a}^ae^{ixy}\varhatf(y)dy
=\sqrt{\Bover2{\pi}}\int_0^a\cos xy\,\varhatf_c(y)dy
+\sqrt{\Bover2{\pi}}\int_0^a\sin xy\,\varhatf_s(y)dy$}

\noindent untuk setiap $x\in\Bbb R$ dan $a\ge 0$.

\spheader 283Xp Gunakan fakta bahwa
$\int_{0}^a\int_{0}^{\infty}e^{-xy}\sin
y\,dxdy=\int_{0}^{\infty}\int_{0}^ae^{-xy}\sin y\,dydx$ setiap kali
$a\ge 0$ untuk menunjukkan bahwa 
$\int_0^{\infty}\bover1{1+x^2}dx=\lim_{a\to\infty}\int_0^a\bover{\sin y}{y}dy$.

\sqheader 283Xq Tunjukkan bahwa jika $f(x)=e^{-\sigma|x|}$, di mana
$\sigma>0$, maka
$\varhatf(y)=\bover{2\sigma}{\sqrt{2\pi}(\sigma^2+y^2)}$.
Oleh karena itu, atau sebaliknya, temukan transformasi Fourier dari
$y\mapsto\Bover{1}{1+y^2}$.

\spheader 283Xr Temukan transformasi Fourier invers dari fungsi karakteristik
dari interval terbatas di $\Bbb R$.   Tunjukkan bahwa
dalam arti formal 283F dapat dianggap sebagai kasus khusus 283O.

\spheader 283Xs Misalkan $f$ adalah fungsi yang dapat diintegralkan nonnegatif pada
$\Bbb R$, dengan transformasi Fourier $\varhatf$.   Tunjukkan bahwa
\discrcenter{468pt}{$\sum_{j=0}^n\sum_{k=0}^n
  a_j\bar a_k\varhatf(y_j-y_k)\ge 0$ }setiap kali $y_0,\ldots,y_n$ di
$\Bbb R$ dan $a_0,\ldots,a_n\in\Bbb C$.

\spheader 283Xt Misalkan $f$ adalah fungsi bernilai kompleks yang dapat diintegralkan pada
$\Bbb R$.   Tunjukkan bahwa $\tilde f(x)=\sum_{n=-\infty}^{\infty}f(x+2\pi n)$
didefinisikan dalam $\Bbb C$ untuk hampir setiap $x$.
\Hint{$\sum_{n=-\infty}^{\infty}\int_{-\pi}^{\pi}|f(x+2\pi n)|dx
<\infty$.}   Tunjukkan bahwa $\tilde f$ bersifat periodik.   Tunjukkan bahwa koefisien Fourier
dari $\tilde f\restr\ocint{-\pi,\pi}$ adalah
$\langle\bover{1}{\sqrt{2\pi}}\varhatf(k)\rangle_{k\in\Bbb Z}$.

\leader{283Y}{Latihan lebih lanjut (a)}
%\spheader 283Ya
Tunjukkan bahwa jika $f:\Bbb R\to\Bbb C$ adalah
kontinu mutlak dalam setiap interval berbatas, $f'$ mempunyai variasi
terbatas pada $\Bbb R$, dan
$\lim_{x\to\infty}f(x)=\lim_{x\to-\infty}f(x)=0$, maka

\Centerline{$g(y)=\Bover1{\sqrt{2\pi}}\lim_{a\to\infty}
\int_{-a}^ae^{-iyx}f(x)dx
=-\Bover{i}{y\sqrt{2\pi}}\lim_{a\to\infty}
\int_{-a}^ae^{-iyx}f'(x)dx$}

\noindent terdefinisi, dengan

\Centerline{$y^2|g(y)|\le\Bover4{\sqrt{2\pi}}\Var_{\Bbb R}(f')$,}

\noindent untuk setiap $y\ne 0$.
%283C

\spheader 283Yb Misalkan $f:\Bbb R\to\Bbb C$ adalah fungsi
yang dapat diintegralkan dan kontinu mutlak pada setiap interval berbatas, dan
anggap bahwa turunannya $f'$ mempunyai variasi terbatas pada $\Bbb R$.
Tunjukkan bahwa $\varhatf$ dapat diintegralkan dan
$f=\varhatf\varcheck{\phantom{f}}$.   \Hint{225Yd, 283Ci, 283Xk.}
%283Xk 283K

\spheader 283Yc Misalkan $f:\Bbb R\to\coint{0,\infty}$ adalah fungsi
genap sehingga $f$ cembung pada $\coint{0,\infty}$
dan $\lim_{x\to\infty}f(x)=0$.

\quad{(i)} Tunjukkan bahwa, untuk $y>0$ dan $k\in\Bbb N$,
$\int_{-2k\pi/y}^{2k\pi/y}e^{-iyx}f(x)dx\ge 0$.

\quad{(ii)} Tunjukkan bahwa
$g(y)=\bover1{\sqrt{2\pi}}\lim_{a\to\infty}\int_{-a}^ae^{-iyx}f(x)dx$
ada di $\coint{0,\infty}$ untuk setiap $y\ne 0$.

\quad{(iii)} Untuk $n\in\Bbb N$, atur $f_n(x)
=e^{-|x|/(n+1)}f(x)$ untuk setiap $x$.   Tunjukkan bahwa $f_n$ dapat diintegralkan dan
cembung di $\coint{0,\infty}$.

\quad{(iv)} Tunjukkan $g(y)=\lim_{n\to\infty}\varhatf_{n}(y)$ untuk
setiap $y\ne 0$.

\quad{(v)} Tunjukkan bahwa jika $f$ dapat diintegralkan maka

\Centerline{$\int_{-a}^a\varhatf
=\Bover4{\sqrt{2\pi}}\int_{0}^{\infty}\Bover{\sin at}{t}f(t)dt
\le\Bover{4a}{\sqrt{2\pi}}\int_0^{\pi/a}f
\le 2\sqrt{2\pi}f(0)
$}

\noindent untuk setiap $a\ge 0$.   Oleh karena itu, tunjukkan bahwa, baik $f$
dapat diintegralkan maupun tidak, $g$ dapat diintegralkan dan
$f_n=(\varhatf_n)\varspcheck$ untuk setiap $n$.

\quad{(vi)} Tunjukkan bahwa
$\lim_{a\downarrow 0}\sup_{n\in\Bbb N}\int_{-a}^a\varhatf_n=0$.

\quad(vii) Tunjukkan bahwa jika $f'$ dibatasi (pada domainnya) maka
$\{\varhatf_n:n\in\Bbb N\}$ dapat diintegralkan secara seragam ({\it petunjuk\/}: gunakan
(vi) dan 283Ya), sehingga $\lim_{n\to\infty}\|\varhat f_n-g\|_1=0$ dan
$f=\varcheck{g}$.

\quad(viii) Tunjukkan bahwa jika $f'$ tidak terikat maka untuk setiap $\epsilon>0$ kita
dapat menemukan $h_1$, $h_2:\Bbb R\to\coint{0,\infty}$, keduanya genap, cembung dan
konvergen ke $0$ di $\infty$, sehingga $f=h_1+h_2$, $h'_1$ terikat,
$\int h_2\le\epsilon$ dan $h_2(0)\le\epsilon$.   Oleh karena itu tunjukkan bahwa dalam kasus
ini juga $f=\varcheck{g}$.

\spheader 283Yd Misalkan $f:\Bbb R\to\Bbb R$ genap, terdiferensiasi dua
kali, dan konvergen ke $0$ pada $\infty$; misalkan $f''$ kontinu dan
$\{x:f''(x)=0\}$ terbatas dalam $\Bbb R$. Tunjukkan bahwa $f$ merupakan
transformasi Fourier dari suatu fungsi yang dapat diintegralkan.
\Hint{Gunakan 283Yc dan 283Yb.}
%283 notes

\spheader 283Ye Misalkan $g:\Bbb R\to\Bbb R$ adalah fungsi ganjil dari variasi terbatas
sedemikian rupa sehingga $\int_1^{\infty}\bover1xg(x)dx=\infty$.
Tunjukkan bahwa $g$ tidak bisa menjadi transformasi Fourier dari fungsi apa pun yang dapat diintegralkan
$f$.   ({\it Petunjuk\/}: tunjukkan jika $g=\varhatf$ maka
%283Yd

\Centerline{$-i\int_0^1f
=\Bover{2}{\sqrt{2\pi}}\lim_{a\to\infty}
  \int_0^a\Bover{1-\cos x}{x}g(x)dx=\infty$.)}
}%end of exercises

\endnotes{
\Notesheader{283}
Dalam bagian ini saya telah mencoba menyajikan teori dasar transformasi
Fourier untuk fungsi yang dapat diintegralkan pada $\Bbb R$, dengan
mempertimbangkan perluasan konsep yang akan dicoba dalam bagian berikutnya.
Seperti dalam \S282, saya menonjolkan dua teorema (283I dan 283L) yang
memberikan syarat agar inversi transformasi Fourier mengembalikan fungsi
semula; kita berhadapan dengan integral tak wajar
$\lim_{a\to\infty}\int_{-a}^a$, sebagaimana sebelumnya kita berhadapan dengan
jumlah simetris $\lim_{n\to\infty}\sum_{k=-n}^n$. Saya tidak melangkah sejauh
di \S282; khususnya, untuk sementara saya menunda pembahasan fungsi
terintegralkan-kuadrat, karena transformasi Fourier-nya mungkin tidak dapat
dinyatakan dengan rumus sederhana yang digunakan di sini.

Salah satu hambatan paling mendasar dalam bidang ini adalah tidak adanya
kriteria efektif untuk menentukan fungsi mana yang merupakan transformasi
Fourier dari fungsi yang dapat diintegralkan. (Untungnya, keadaan lebih baik
untuk fungsi terintegralkan-kuadrat; lihat 284O-284P.) Dalam 283Yc-283Yd saya
meringkas argumen yang menunjukkan bahwa fungsi {\it genap} `biasa' yang tidak
berosilasi dan konvergen ke $0$ pada $\pm\infty$ merupakan transformasi Fourier
dari fungsi yang dapat diintegralkan. Yang mencolok, hal ini tidak berlaku bagi
fungsi {\it ganjil}; jadi $y\mapsto\Bover{1}{\ln(e+y^2)}$ merupakan transformasi
Fourier dari suatu fungsi yang dapat diintegralkan, tetapi
$y\mapsto\Bover{\arctan y}{\ln(e+y^2)}$ bukan merupakan contoh (283Ye).

Di 283W saya membuat sketsa teori transformasi Fourier yang sesuai
$\Bbb R^r$.   Ada sedikit kejutan.   Satu hal yang perlu diperhatikan adalah di mana
dalam kasus satu dimensi kita meminta turunan kedua yang berkelakuan baik,
dalam kasus dimensi $r$ kita mungkin perlu mendiferensiasikan $r+1$ kali
(283Wf).   Alasan lainnya adalah kita kehilangan `prinsip lokalisasi'.   Di dalam
kasus satu dimensi, jika $f$ dapat diintegralkan dan nol pada suatu interval
$\ooint{c,d}$, lalu
$\lim_{a\to\infty}\int_{-a}^ae^{ixy}\varhatf(y)dy=0$ untuk setiap
$x\in\ooint{c,d}$;  ini langsung dari 283I atau 283L.   Namun dalam
dimensi yang lebih tinggi formulasi paling alami dari hasil yang sesuai
tidak berlaku (283Wm).




}%end of notes

\discrpage
